\begin{abstract}
Test-time augmentation (TTA) is widely treated as a low-risk way to improve model accuracy at inference time, yet a growing body of evidence shows it can also hurt. We conduct a preregistered, statistically audited study of naive mixed-policy TTA on three MedMNIST datasets (PathMNIST, BloodMNIST, DermaMNIST) across two normalization schemes (BatchNorm, GroupNorm), two resolutions (28px, 64px), and three training seeds, plus a fixed 128px extension and a DermaMNIST/ResNet-18 positive-control replication. Within 30 distinct trained-model cells evaluated under a fixed mixed augmentation policy at N=50 views, naive TTA reduced test accuracy relative to a single clean forward pass in every cell, with per-cell drops ranging from 18.76 to 66.09 percentage points; these 30 cells span roughly eight to ten underlying dataset/resolution/normalization configurations, each replicated over three seeds, and the claim is a within-model one rather than a population replication. A separate, non-preregistered secondary analysis comparing models trained with a training-time augmentation policy matched to the test-time TTA policy against models trained without that match found all six fixed-model paired comparisons favoring the matched-policy models, each 95\% confidence interval excluding zero. Normalization and resolution showed heterogeneous, dataset-dependent secondary patterns rather than a single consistent mitigation. A positive-control replication (DermaMNIST, ResNet-18) did not reproduce the positive TTA effect reported by the source study we build on. A post-review extension shows the harm is stable across view counts from 1 to 100 (worst at N=1); that the labeled diagnostic structure survives in the large majority of augmented views (human-judged lost in 0-4\% per dataset); and that a geometric-only augmentation family --- which preserves image content --- still causes harm on its own, indicating the effect is a failure of augmentation robustness rather than label destruction. We report these findings using an explicit two-tier evidentiary framework --- preregistered within-cell evidence versus secondary fixed-model comparisons --- to avoid overstating what a non-preregistered cross-condition comparison can support, and we release a fully audited, hash-verified evidence and reproducibility trail alongside this manuscript.
\end{abstract}

\section{Introduction}

Test-time augmentation applies a set of input transformations at inference time and aggregates the resulting predictions, most commonly by averaging predicted class probabilities. It is often presented as a nearly-free accuracy improvement, since it requires no retraining and only additional forward passes. \citet{medeiros2026tta} recently challenged this assumption directly, reporting that mixed-policy TTA with mean aggregation at 50 views degraded accuracy in 11 of 12 model-dataset combinations tested on three MedMNIST benchmarks, with drops as large as 31.6 percentage points and a single improving case (ResNet-18 on DermaMNIST, +1.6 percentage points). This result runs against the common practitioner intuition that TTA is a safe default, and --- because it is a single-author preprint with no peer-reviewed venue found at the time of our review --- warrants independent, methodologically strengthened verification before it is treated as established.

This project undertakes that verification. We do not treat the source preprint's numbers as settled fact; we design an independent, preregistered confirmatory matrix over the same three MedMNIST datasets, add a normalization ablation (BatchNorm vs. GroupNorm) and a resolution ablation (28px vs. 64px, with a 128px extension), add a training-policy-matching arm, and replicate the source study's single reported positive result (DermaMNIST, ResNet-18) as an internal positive control. Every hypothesis, endpoint, and statistical procedure was frozen before any test-split evaluation occurred, and the frozen statistical analysis was executed by a sealed, cryptographically fingerprinted pipeline rather than run ad hoc after seeing results (see "Reproducibility and Audit Trail").

Our central empirical finding replicates and generalizes the direction of the source preprint's headline result: naive TTA harmed every one of 30 distinct trained-model cells we evaluated under the fixed mixed- policy/N=50 condition. Beyond replication, we ask whether this harm is mitigated by training-time exposure to the same augmentation distribution used at test time, and whether normalization or resolution choices modulate the harm. We report the answers to these questions using an explicit two-tier evidentiary framework, because the normalization, resolution, and policy-matching questions were only formally specified as fixed-model comparisons \emph{after} our validation- stage results were already observed and \emph{before} the test split was opened --- a design choice we disclose plainly rather than presenting these comparisons with the same confirmatory status as our preregistered within-cell tests. We also report, without qualification softening the result, that our internal positive-control replication (DermaMNIST, ResNet-18) did not reproduce the source preprint's reported positive TTA effect.

This manuscript makes no claim of state-of-the-art performance, no clinical-utility or diagnostic claim (MedMNIST is explicitly not validated for clinical use), and no claim that our findings generalize to medical imaging models broadly. Its contribution is a controlled, statistically disciplined, fully audited account of when and how far naive TTA harm extends within the specific models, datasets, resolutions, and normalization choices we actually evaluated --- and an explicit demonstration of how much confidence a non-preregistered secondary comparison can and cannot support.

The motivation for this level of methodological care is practical, not rhetorical. A practitioner deciding whether to enable TTA in a deployed image-classification pipeline typically has access to exactly the kind of evidence this study distinguishes carefully: a small number of validation-set observations, and a temptation to compare conditions that were never designed as a controlled experiment. If a $-$60 percentage-point accuracy collapse under naive TTA is possible on ordinary, non-adversarial MedMNIST inputs --- as our within-cell results below show it is, in this matrix --- then treating TTA as a low-risk default is unsafe without validating it on the specific model and augmentation policy in use. Conversely, if a secondary, non-preregistered comparison is reported with the same confidence as a preregistered one, a practitioner may draw a stronger conclusion (e.g., "GroupNorm always mitigates TTA harm") than the evidence supports, since our own normalization result reverses sign by dataset. We therefore treat the preregistered-versus-secondary distinction not as a statistical formality but as the paper's main methodological contribution alongside its empirical findings: every numeric claim below is labeled with the tier of evidence it is drawn from, and Table 1 in the Experimental Design section makes that labeling explicit before any result is presented.

\textbf{Contributions.} (i) A preregistered, statistically audited replication showing naive mixed-policy TTA reduced test accuracy in every one of 30 distinct trained-model cells on three MedMNIST datasets, with a fully hash-verified, independently re-checkable evidence chain. (ii) An end-to-end, mechanically enforced separation of preregistered within-cell evidence from secondary cross-condition evidence, offered as a reusable reporting discipline for empirical TTA studies. (iii) Secondary fixed-model evidence that training-time / test-time augmentation-policy matching mitigates the harm, reported with its non-confirmatory status made explicit. (iv) A post-review mechanistic extension --- a view-count scaling curve, a label- preservation audit, and a per-augmentation-component decomposition --- that localizes the effect as a failure of augmentation robustness (present even for a content-preserving geometric-only policy) rather than a consequence of the augmentation destroying labeled content. (v) A public release of the full evidence, audit trail, and a standalone one-command reproduction.

\section{Related Work}

\textbf{TTA can flip correct predictions to incorrect even when net accuracy improves.} \citet{shanmugam2021better} show that naive averaging across augmented views in TTA can change many individual correct predictions to incorrect ones even in settings where the net accuracy effect is positive, and propose a learned aggregation-weighting scheme as a fix. This "harm rate" framing --- that aggregate accuracy can hide sample-level harm --- motivates our own descriptive seed-level reporting discipline, though we do not adopt their learned aggregation method; we cite this work as the origin of the harm-framing concept rather than presenting it as our own.

\textbf{Validation-driven and per-sample learned TTA policy selection already exist.} \citet{lyzhov2020greedy} introduce Greedy Policy Search (GPS), a non-learned-network method that greedily selects a test-time augmentation policy to maximize validation performance. \citet{kim2020learning} go further, learning an auxiliary network that predicts per-transform loss for each test input and filters transforms per sample before aggregation, i.e. instance-level selective TTA. Neither of these compares a matched-versus-unmatched training-time augmentation policy as a fixed- model secondary analysis, which is the specific angle of our own policy-matching arm; we note this as the narrowest gap we identified, not as an exhaustively verified absence, since we did not perform a systematic citation-graph search.

\textbf{TTA aggregation has been evaluated specifically in medical imaging.} \citet{sherkatghanad2024baytta} propose BayTTA, a Bayesian model averaging scheme for combining TTA-view predictions, evaluated on skin cancer, breast cancer, and chest X-ray classification tasks. This is direct prior art for TTA aggregation methods evaluated in a medical-imaging context; our study does not claim to be the first to study TTA in medical imaging, and does not evaluate Bayesian aggregation. \citet{disalvo2024medmnistc} build MedMNIST-C, a corruption-robustness benchmark on top of MedMNIST+, and show that their corruption transforms improve robustness when used as \emph{training-time} data augmentation --- a different use of "augmentation" than the test-time aggregation harm we study, but directly relevant background on the MedMNIST+ resolution family we also use.

\textbf{BatchNorm's running statistics are a documented distribution-shift liability.} \citet{schneider2020improving} show that recomputing or adapting BatchNorm activation statistics at test time using unlabeled test-distribution statistics, instead of relying on training-time running statistics, produces large corruption-robustness gains on ImageNet-C. This supports the plausibility of BatchNorm's running statistics as one mechanism by which augmented test inputs could degrade accuracy, and is background for our normalization ablation; we do not adapt BatchNorm statistics ourselves and make no claim about the specific mechanism underlying our observed harm beyond what is stated as a hypothesis in the Discussion. Group Normalization \citep{wu2018group} is the batch-independent alternative we compare against BatchNorm: it normalizes within per-sample channel groups rather than relying on batch or running statistics. We note explicitly that GroupNorm differs from BatchNorm in more properties than the mere presence of running statistics, so any normalization-conditioned difference we observe is evidence \emph{consistent with}, not \emph{proof of}, a running-statistics mechanism.

\textbf{Test-time adaptation, medical-imaging TTA for uncertainty, and theory of TTA are adjacent lines we do not extend.} Beyond static BatchNorm-statistics recomputation, TENT \citep{wang2021tent} adapts a model online at test time by minimizing prediction entropy over each test batch, updating BatchNorm's affine parameters --- a more active form of test-time adaptation than anything we evaluate; our study fixes the trained model and only aggregates augmented-view predictions. In medical imaging specifically, \citet{ayhan2018ttaug} use test-time augmentation not to improve point accuracy but to \emph{estimate} predictive uncertainty from the variability of predictions across augmented views; our harm/rescue-rate reporting is adjacent to this uncertainty framing but we do not compute calibration-as-uncertainty from augmented-view dispersion. On the theory side, \citet{kimura2024understanding} analyzes when TTA can help or hurt in terms of the diversity and error-correlation of the augmented-view predictions; this is a complementary lens to the distribution-shift account we sketch in the Discussion, and we do not attempt a formal treatment. We list these to scope our contribution: we study \emph{naive} mean-probability TTA with a fixed model, and do not propose or evaluate test-time adaptation, an uncertainty estimator, or a theoretical bound.

\textbf{MedMNIST is the dataset family underlying this study.} MedMNIST v2 (Yang et al., 2023, published in \emph{Scientific Data}) provides standardized 28px classification benchmarks across biomedical imaging modalities, explicitly stating the datasets are \textbf{not intended for clinical use} \citep{yang2023medmnistv2}. The earlier MedMNIST Classification Decathlon (Yang et al., 2021, ISBI) introduced the original 28px benchmark suite \citep{yang2021medmnist}. MedMNIST+ extends select subsets to 64px, 128px, and 224px, constructed from independently higher-resolution source images rather than upsampled from the 28px files, preserving sample indices and splits across resolutions --- a fact we rely on directly for our resolution ablation's validity (see Methods).

\textbf{Corruption-robustness benchmarking on MedMNIST is a separate, active line of work.} MedMNIST-C's corruption taxonomy (blur, noise, digital, and weather-style corruptions adapted to the 28px-family MedMNIST images) targets robustness to naturally-occurring image corruption, and its reported improvements come from \emph{training} models on corrupted data, not from aggregating predictions over augmented views at test time \citep{disalvo2024medmnistc}. We cite it here as evidence that the MedMNIST+ resolution family we use is already an active substrate for augmentation-related robustness research, and as a candidate source of a validated corruption/augmentation taxonomy for future extensions of this work, not as directly overlapping prior art for the TTA-harm question itself.

\textbf{Novelty scope.} Taken together, the reviewed literature establishes that TTA harm, augmentation-policy mismatch as a plausible mitigation lever, medical-imaging-specific TTA evaluation, test-time adaptation, and a theory of when TTA helps are all active, documented areas, and this study does not claim to originate, or to advance the method frontier of, any of them. We propose no new TTA method, aggregator, adaptation procedure, uncertainty estimator, benchmark, or theoretical bound. What this study offers is a preregistered, fully audited replication and characterization: an independent verification of a circulating preprint's headline claim, a mechanically enforced preregistered-versus-secondary reporting discipline, and a post-review mechanistic localization (scaling curve, label-preservation audit, per-component decomposition) showing the harm is an augmentation-robustness effect rather than label destruction. The specific empirical combination we did not find directly addressed in the sources reviewed --- matched-versus-unmatched training policy as a fixed-model secondary difference-in-differences analysis, alongside a controlled normalization/resolution ablation, on the MedMNIST family, under a two-tier reporting discipline --- is offered as a narrow gap from a non-exhaustive review (no citation-graph traversal), not a claimed first.

\section{Methods}

\textbf{Datasets.} We use three MedMNIST v2 / MedMNIST+ subsets \citep{yang2023medmnistv2,yang2021medmnist}: PathMNIST (train/val/test = 89,996 / 10,004 / 7,180; CC BY 4.0), BloodMNIST (11,959 / 1,712 / 3,421; CC BY 4.0), and DermaMNIST (7,007 / 1,003 / 2,005; CC BY-NC 4.0, hence non-commercial). Higher-resolution variants (64px, 128px) are constructed by the MedMNIST+ project from independently-sourced higher-resolution originals --- PathMNIST and DermaMNIST by direct resizing (from 224x224 and 600x450 sources respectively, no cropping), BloodMNIST by center-cropping a 360x363 source to 200x200 before resizing --- with sample indices and train/val/test splits preserved identically across resolutions. This means our resolution ablation (H2) measures the effect of retained source-image information at a given standardized resolution, not an artifact of upsampling the 28px images.

\textbf{Models.} A small paper-constrained CNN ("SmallCNN," 94,857 parameters, three convolutional layers with global average pooling so parameter count is resolution-invariant, evaluated with both BatchNorm and GroupNorm normalization), and a ResNet-18 variant adapted for 28x28-family inputs (approximately 11M parameters, BatchNorm only, no initial pooling layer), following the source study's architecture choices where they are specified there.

\textbf{Training.} All confirmatory models are trained with Adam (learning rate 0.001, no weight decay), cross-entropy loss, up to 30 epochs with cosine-annealed learning rate, early stopping on validation accuracy (patience 5 epochs, no minimum-improvement threshold), restoring the best-validation-accuracy checkpoint, batch size 256, float32 precision, no label smoothing or class weighting, and inputs scaled from uint8 to [0,1] with no additional channel standardization. Training-time augmentation is absent everywhere except the policy-matching arm (see below). These operational choices (exact epoch cap, patience, weight-decay, and minimum-improvement values) are preregistered choices of this study, not values extracted from the source preprint, which reports only "25-30 epochs" and does not specify the others; this remains a paper-constrained reproduction, not an exact replication.

\textbf{TTA procedure.} Test-time augmentation applies a fixed \emph{mixed} policy (geometric: horizontal/vertical flips, +/-15 degree rotation, random resized crop 0.8-1.0x; intensity: brightness/contrast jitter +/-0.3, Gaussian blur; mixed: both families combined) over a deterministic, per-sample registered 100-view sequence, of which the primary confirmatory condition uses the first 50 views (N=50), aggregated by mean predicted probability --- reproducing the source study's headline condition.

\textbf{Policy-matching arm.} A subset of models is additionally trained with the same mixed augmentation policy applied once per training sample per training step, so that the training-time and test-time augmentation distributions match. This is compared against otherwise- identical models trained without that augmentation (the naive-TTA condition above) to test whether policy matching mitigates TTA harm.

\textbf{Data licensing and preprocessing.} PathMNIST and BloodMNIST are distributed under CC BY 4.0; DermaMNIST is distributed under CC BY-NC 4.0, and we accordingly treat every DermaMNIST-derived number in this manuscript as non-commercial-use-only. No dataset-specific channel mean or standard deviation is applied; inputs are scaled from their native uint8 range to [0, 1] and otherwise used as distributed. We do not resize, crop, or otherwise modify the MedMNIST+-provided resolution variants ourselves --- the 28px, 64px, and 128px inputs are used exactly as constructed by the MedMNIST+ project, so any resolution effect we observe is attributable to the MedMNIST+ construction process described above, not to a resizing choice of our own.

\textbf{Compute environment.} All training and evaluation was executed on a single Apple Silicon machine (Apple M3 Pro, 18GB unified memory) using the PyTorch MPS backend, float32 precision throughout, with no CUDA hardware involved. This is disclosed because it constrains the practical scale of the confirmatory matrix (see Experimental Design and Limitations) relative to what a larger GPU cluster would allow, and because MPS-backend numerics can differ in minor ways from CUDA reference implementations; we did not cross-validate our numbers against a CUDA run.

\section{Experimental Design}

The confirmatory matrix comprises four blocks. \textbf{Block A} (24 runs) crosses PathMNIST/BloodMNIST, 28px/64px, SmallCNN with BatchNorm/ GroupNorm, and three seeds, isolating the normalization (H1) and resolution (H2) comparisons under a fixed unmatched TTA policy. \textbf{Block B} (6 runs) trains SmallCNN/BatchNorm at 28px on PathMNIST and BloodMNIST with the training-time augmentation policy matched to the test-time TTA policy, across three seeds, for the policy-matching comparison (H3). \textbf{Block C} (3 runs) trains ResNet-18/BatchNorm on DermaMNIST at 28px across three seeds as an internal positive-control replication of the source study's sole reported improving case. \textbf{Block D} (6 runs, conditional on a compute-budget gate) extends the normalization/resolution matrix to a fixed 128px resolution for PathMNIST and BloodMNIST with BatchNorm. Every training seed (0, 1, 2) is drawn from three fixed confirmatory seeds; the exploratory pilot seed used in earlier, non-confirmatory pipeline validation work is permanently excluded from every confirmatory run.

Every model is evaluated under a clean (single, unaugmented forward pass) condition and the naive-TTA (N=50, mixed policy, mean probability) condition described above. This produces, across all executed blocks, \textbf{39 total preregistered matrix cells}. Of these, 30 are distinct trained models evaluated under the unmatched (naive-TTA) policy: Block A's 24 cells, Block D's 6 cells (Block D is resolution/normalization-only and therefore contributes only to the unmatched arm), and Block B's 6 cells contribute their own 6 unmatched- arm comparisons but those 6 run\_ids coincide with 6 of Block A's run\_ids at the shared dataset/resolution/normalization/seed identity, so they are \textbf{not} additional distinct cells --- overlapping hypothesis- family membership (a given trained model can be part of more than one hypothesis's evaluated cell set) never inflates the distinct-cell count. The 30-cell figure used throughout this manuscript is the deduplicated count.

\textbf{Block D's conditional status.} Because 128px training and evaluation is substantially more compute-intensive than 28px/64px on our single- machine setup, Block D was preregistered as \emph{approved-conditional}: it executes only if a native 128px runtime/memory benchmark (measured before any 128px training) confirms the block fits the project's compute budget. This conditional-execution rule was frozen before any 128px result existed, so Block D's inclusion in the final 39-cell matrix is not a post-hoc convenience --- it was always a possible outcome of the frozen protocol, and its 6 cells are evaluated identically to Block A's cells once activated.

\textbf{Evidentiary tiers.} We distinguish two categories of comparison throughout. \textbf{Preregistered within-cell comparisons} (clean vs. naive- TTA accuracy on the \emph{same} trained model) were specified before any test-split evaluation and constitute our confirmatory evidence. \textbf{The normalization (H1), resolution (H2), and policy-matching (H3) cross-condition comparisons} --- i.e., comparing \emph{different} trained models' TTA effects against each other via a fixed-model difference-in-differences (DiD) estimate --- were specified as a fixed statistical procedure after our validation-stage results were already observed, but before the test split was opened. This ordering means these three comparisons cannot carry the same preregistered-confirmatory status as the within-cell tests; we report them as secondary, post-validation/pre-test-specified evidence throughout, and we do not apply any pooled or population-level significance procedure across either tier.

\begin{table}[t]
\centering
\small
\begin{tabular}{@{}p{2.6cm}p{3.4cm}p{2.4cm}p{4.4cm}@{}}
\toprule
\textbf{Evidence tier} & \textbf{Source} & \textbf{Cells / pairs} & \textbf{Confirmatory?} \\
\midrule
Preregistered within-cell & H1/H2/H3/BLOCK\_C & 39 unique cells & Yes --- clean-vs-TTA, within one trained model \\
Secondary fixed-model comparison & Cross-condition H1/H2/H3 & 30 pairs (12+12+6) & No --- post-validation / pre-test-specified, never preregistered \\
Descriptive summary & Preregistered seed-level groupings & 13 dataset $\times$ resolution $\times$ normalization groups & No --- non-inferential, no p-value / CI of its own \\
Unsupported / forbidden & H4; pooled or model-population verdicts; secondary significance labels & N/A & Never permitted anywhere in this package \\
\bottomrule
\end{tabular}
\caption{\textbf{Experimental-design and evidence-classification map}, referenced throughout Results and Discussion. Reproduced from the committed evidence package.}
\label{tab:design}
\end{table}

\textbf{Table 1} (reproduced from the committed evidence package, \texttt{artifacts/\allowbreak paper\_\allowbreak evidence/\allowbreak tables/\allowbreak table\_\allowbreak 1\_\allowbreak design\_\allowbreak classification.md}) makes this classification explicit before any result is presented: preregistered within-cell evidence (H1/H2/H3/BLOCK\_C, 39 unique cells) is confirmatory; secondary fixed-model comparisons (cross-condition H1/H2/H3, 30 pairs across the three families) are non-confirmatory; descriptive seed-level summaries carry no inferential value of their own; and a fourth row lists what is unsupported/forbidden at every tier --- H4, any pooled or model-population verdict, and any secondary significance language. We refer back to this table's tier labels throughout the Results and Discussion sections rather than restating the tier-boundary reasoning at every mention.

\section{Statistical Analysis}

For every preregistered within-cell comparison, we compute a paired bootstrap 95\% confidence interval on the delta accuracy (TTA minus clean, on the same test samples from the same trained model), using 10,000 resamples with replacement over the test set, together with an exact or continuity-corrected McNemar test on the paired correct/incorrect 2x2 contingency table --- the appropriate paired binary-outcome test for this design. Within each hypothesis family (H1, H2, H3, and the BLOCK\_C positive control), Benjamini-Hochberg false-discovery-rate correction is applied across that family's cells; both raw and BH-adjusted p-values are reported. A cell that is a member of more than one hypothesis family (see Experimental Design) therefore has one raw McNemar p-value --- the same underlying paired computation --- but a \emph{different} BH-adjusted p-value in each family's own correction set, since each family's FDR correction is computed independently over its own cell membership; we report this per-family rather than collapsing it to a single implied value.

For the three secondary cross-condition comparisons (H1 normalization, H2 resolution, H3 policy-matching), we compute a paired-bootstrap 95\% confidence interval on the difference-in-differences estimate between two fixed, already-trained models' TTA effects, again using 10,000 resamples. We report the point estimate and interval and describe, factually, whether the interval excludes zero. We do not compute, and do not report, a p-value, an alpha threshold, or a significance decision for any secondary comparison, and we never use the word "significant" to describe one, since these comparisons were not preregistered and are not corrected as part of any confirmatory family.

Descriptive seed-level summaries (mean, sample standard deviation, min, max of the per-seed delta-accuracy point estimate within a dataset x resolution x normalization group) are reported separately, carry no p-value or confidence interval of their own, and are labeled non-inferential throughout; they contextualize seed-to-seed variability and are never used as evidence for or against any hypothesis.

Alongside delta accuracy and its confidence interval, effect sizes are reported for each preregistered cell: the raw delta-accuracy magnitude itself, and harm/rescue rates (the fraction of samples correctly classified clean but incorrectly classified under TTA, and vice versa) as effect-size-like quantities in the spirit of \citet{shanmugam2021better}'s observation that net accuracy can mask sample-level flips. These are descriptive characterizations of a single preregistered comparison, not additional hypothesis tests, and do not receive their own confidence interval or p-value.

Pairing is preserved throughout: because clean and TTA predictions are always evaluated on the same test samples from the same trained model, every within-cell comparison is a paired design, and every cross-condition secondary comparison compares two \emph{already independently paired} delta-accuracy distributions from two different fixed models --- never an independent-groups test applied to what is actually paired data. The choice of Benjamini-Hochberg FDR correction per hypothesis family, rather than a single Bonferroni correction pooled across all four families, follows the frozen statistical analysis plan's reasoning that the four hypotheses ask distinct scientific questions and a single pooled correction would be needlessly conservative without a corresponding scientific justification; this choice was fixed before any test-split result existed, not selected after seeing which correction produced a more favorable outcome.

No new statistic beyond the above was computed for this manuscript. Every number reported in Results below is either copied directly from the canonical, hash-verified evidence package described in "Reproducibility and Audit Trail," or is a minimum/maximum taken directly across an already-computed set of per-cell point estimates for narrative purposes (e.g., the range of the 30 within-cell deltas); no new inferential test, pooled estimate, or aggregate hypothesis test was introduced beyond what the frozen statistical analysis plan already specifies.

\section{Results}

\begin{figure*}[t]
\centering
\includegraphics[width=0.95\textwidth]{figure_1_unmatched_policy_forest}
\caption{\textbf{Preregistered within-cell evidence.} Per-cell $\Delta$ accuracy (TTA $-$ clean) with paired-bootstrap 95\% CIs for the 30 distinct unmatched-policy cells at $N{=}50$. Every interval lies below zero.}
\label{fig:unmatched}
\end{figure*}

\textbf{All 30 distinct unmatched-policy cells showed harm.} Under the fixed mixed-policy, N=50, mean-aggregation naive-TTA condition, every one of the 30 distinct trained-model cells evaluated showed a negative delta accuracy (TTA minus clean), ranging from -18.76 percentage points (PathMNIST, 28px, BatchNorm, seed 0) to -66.09 percentage points (BloodMNIST, 64px, BatchNorm, seed 1) (Figure 1; full per-cell values in Supplementary Table 2). Split by dataset, the 15 BloodMNIST cells ranged from -20.70 to -66.09 percentage points and the 15 PathMNIST cells ranged from -18.76 to -51.92 percentage points --- both datasets show substantial harm across their full cell sets, with no cell in either dataset showing a non-negative delta. This is a preregistered, within-cell result: it establishes that naive TTA harmed these specific trained models under this specific policy and view budget, not a claim about medical imaging models in general. Sample sizes underlying each cell's paired bootstrap and McNemar test equal the corresponding dataset's official test-split size (PathMNIST 7,180; BloodMNIST 3,421), so the harm reported here is measured against the full official test set of each dataset, not a subsample.

\begin{figure}[t]
\centering
\includegraphics[width=0.72\textwidth]{figure_2_matched_policy_mitigation}
\caption{\textbf{Secondary fixed-model evidence (non-confirmatory).} Difference-in-differences of the TTA effect, matched- minus unmatched-policy model, same dataset/seed; all six intervals exclude zero.}
\label{fig:matched}
\end{figure}

\textbf{Matched-policy mitigation is supported by secondary evidence and descriptively corroborated within cells.} The six matched-policy within-cell deltas (Block B, PathMNIST and BloodMNIST, 28px, BatchNorm) were small and mixed in sign (-1.21 to +4.09 percentage points), in sharp contrast to the strongly negative unmatched deltas for the corresponding dataset/resolution/normalization cells --- a descriptive pattern, not a second confirmatory test of the same claim. Separately, the six secondary fixed-model DiD comparisons (matched-policy model minus unmatched-policy model's TTA effect, same dataset/seed) were all positive, ranging from +17.55 to +49.58 percentage points, and \textbf{all six 95\% confidence intervals excluded zero} (Figure 2; Supplementary Table 3). We describe this as secondary evidence supporting mitigation, descriptively corroborated by the separate within-cell pattern --- not as a preregistered cross-condition test.

\begin{figure}[t]
\centering
\includegraphics[width=0.72\textwidth]{figure_3_normalization_heterogeneity}
\caption{\textbf{Secondary, non-confirmatory.} GroupNorm-minus-BatchNorm DiD by dataset; the sign reverses between PathMNIST and BloodMNIST.}
\label{fig:norm}
\end{figure}

\textbf{Normalization shows a dataset-dependent secondary pattern.} Across the 12 secondary H1 (GroupNorm minus BatchNorm) DiD comparisons, 12 of 12 95\% confidence intervals excluded zero, but the direction reversed by dataset: 5 of 6 BloodMNIST pairs were positive (GroupNorm showing less harm than BatchNorm), with 1 of 6 negative; all 6 PathMNIST pairs were negative (BatchNorm showing less harm than GroupNorm) (Figure 3; Supplementary Table 4). Both normalization types experienced harm at every within-cell comparison (Results, first paragraph); their relative difference is dataset-dependent secondary evidence only, and must not be read as a general BatchNorm-vs-GroupNorm verdict.

\begin{figure*}[t]
\centering
\includegraphics[width=0.95\textwidth]{figure_4_resolution_comparison}
\caption{\textbf{Secondary, non-confirmatory.} 64px-minus-28px DiD by dataset; higher resolution does not consistently reduce harm.}
\label{fig:res}
\end{figure*}

\textbf{Resolution did not consistently reduce harm.} Across the 12 secondary H2 (64px minus 28px) DiD comparisons, 9 of 12 95\% confidence intervals excluded zero: the BloodMNIST group trended predominantly negative (4 of 6 pairs negative and excluding zero, indicating \emph{more} harm at 64px than 28px, contrary to the hypothesized mitigating direction), while the PathMNIST group was mixed and closer to null (1 of 6 pairs positive and excluding zero, 1 negative and excluding zero, 4 including zero) (Figure 4; Supplementary Table 5). Neither pattern is a preregistered or confirmatory test of a resolution effect, and higher resolution did not consistently reduce TTA harm in either dataset.

\begin{figure}[t]
\centering
\includegraphics[width=0.72\textwidth]{figure_5_block_c_positive_control}
\caption{\textbf{Preregistered positive control.} DermaMNIST / ResNet-18 per-seed $\Delta$ accuracy; the source study's reported $+1.6$pp improvement did not reproduce.}
\label{fig:blockc}
\end{figure}

\textbf{BLOCK\_C did not reproduce the expected positive TTA effect.} The source study's own reported positive TTA result was a +1.6 percentage point improvement for ResNet-18 on DermaMNIST at N=50 \citep{medeiros2026tta}. Across our three positive-control seeds, delta accuracy was +0.25, -2.89, and +0.05 percentage points; only the seed-1 interval (-4.04, -1.80) excluded zero, and it excluded zero on the \emph{negative} side (Figure 5; Supplementary Table 6). This preregistered within-cell positive-control result did not reproduce the source study's expected positive effect in any of the three seeds, and we report this plainly as a null/contrary finding rather than omitting or softening it.

\setcounter{table}{6}
\begin{table}[t]
\centering
\small
\begin{tabular}{@{}p{4.6cm}p{4.0cm}p{4.4cm}@{}}
\toprule
\textbf{Claim} & \textbf{Evidence tier} & \textbf{Status} \\
\midrule
Naive TTA harmed all 30 distinct unmatched-policy base cells & Preregistered within-cell & Supported \\
Matched-policy training mitigates TTA harm & Secondary fixed-model DiD, descriptively corroborated by separate within-cell patterns & Supported only secondarily / descriptively --- not a preregistered cross-condition test \\
Normalization changes the magnitude of harm & Secondary fixed-model DiD only & Supported only secondarily; direction is dataset-dependent \\
Higher resolution reduces TTA harm & Secondary fixed-model DiD only & Contradicted for BloodMNIST; mixed / near-null for PathMNIST \\
BLOCK\_C reproduces the source paper's positive TTA improvement & Preregistered within-cell (positive control) & Contradicted --- expected improvement not reproduced in any seed \\
Any H4 (Validation-Gated TTA) verdict & None --- no derivable family exists & Not made, anywhere \\
Any model-population or general medical-imaging generalization & None --- three seeds, fixed policy / budget & Not permitted, anywhere \\
\bottomrule
\end{tabular}
\caption{\textbf{Claim adjudication.} Each headline statement, the tier of evidence behind it, and its status. Reproduced from the committed evidence package.}
\label{tab:claims}
\end{table}

\textbf{Summary of claim status.} Table 7 (claim adjudication) summarizes which of the claims above are preregistered/confirmatory versus secondary/descriptive, and explicitly lists claims this study does not make. No H4 claim is made anywhere in this manuscript: Validation-Gated TTA (H4) was never implemented and has no derivable evidence family in the canonical evidence package. No population-level or model-population inference is made anywhere in this manuscript, no secondary significance claim is made, and no claim is made that BLOCK\_C reproduced the external reference.

\section{Discussion}

Our central, preregistered finding --- that naive mixed-policy TTA harmed accuracy in all 30 distinct trained-model cells we evaluated --- replicates the direction of Medeiros (2026)'s headline result on an independent, statistically audited implementation, and extends it with a normalization and resolution ablation the source study did not report. That every cell we tested showed harm, across two architectures, two normalization schemes, two resolutions, and three MedMNIST datasets, is a stronger and more systematic demonstration of TTA harm's breadth within this specific matrix than a single-condition result would be --- but it remains a within-model claim, not a claim about TTA or medical imaging models generally, and it says nothing about model-dataset combinations, architectures, augmentation policies, or view budgets we did not test.

The secondary policy-matching result is, in our view, the most practically suggestive finding, precisely because it is also the one we are most careful to qualify. All six fixed-model comparisons favored models whose training-time augmentation distribution matched the test-time TTA policy, and this pattern is descriptively consistent with the separate within-cell observation that matched-policy models showed near-zero or mixed TTA effects while unmatched-policy models showed large, uniformly negative effects. One plausible explanation --- offered here explicitly as a \textbf{hypothesis for future work, not an established mechanism} --- is that naive TTA harm is substantially a distribution- shift effect: augmented test views resemble a distribution the unmatched model never saw during training, while a matched-policy model has already adapted to that distribution. This is consistent with, but not proven by, our data, and is also consistent with the BatchNorm- running-statistics mechanism \citet{schneider2020improving} study in a different context; we did not run the adaptation or ablation experiments (e.g., recomputing BatchNorm statistics at test time, or a \texttt{no\_\allowbreak running\_\allowbreak stats} decomposition arm) that would be needed to isolate which mechanism, or combination of mechanisms, is responsible.

The normalization and resolution results are, by contrast, genuinely heterogeneous rather than a story we can simplify. If BatchNorm's running-statistics dependence were the whole story, we would expect GroupNorm to consistently show less harm than BatchNorm; instead the direction reversed by dataset. If higher resolution reliably preserved more discriminative signal against augmentation-induced distortion, we would expect 64px to consistently show less harm than 28px; instead BloodMNIST trended in the opposite direction and PathMNIST was largely null. We report both patterns as dataset-dependent secondary findings precisely because they resist a single mechanistic narrative, and we consider this heterogeneity itself an informative (if not confirmatory) result: it is evidence against any simple, dataset-independent normalization or resolution rule for predicting or mitigating TTA harm within this matrix.

Our positive-control replication failure (BLOCK\_C) is a genuine null result relative to the source study's own reported case, not a methodological artifact we can explain away; we report it as such. It does not, on its own, invalidate the source study's finding (a single non-reproduction across three seeds, on a positive effect the source study itself reported as small (+1.6 percentage points) relative to its other results, is limited evidence), but it does mean our own dataset does not provide independent confirmation of that specific positive case, and we do not claim otherwise.

Finally, we note that the magnitude of harm we observe in our preregistered cells (up to -66.09 percentage points) exceeds the largest drop reported by the source study (-31.6 percentage points for ResNet-18 on PathMNIST) \citep{medeiros2026tta}. We do not interpret this as evidence that our implementation is more "correct" or that the source study understated the effect; the two studies differ in exact architecture instantiation, training hyperparameters (several of which the source study does not fully specify, as noted in Methods), TTA view sequencing, and dataset/resolution combinations tested, any of which could account for the difference. We report the comparison only to situate our numbers relative to the literature, not to adjudicate whose measurement is more accurate --- a claim neither this manuscript nor a single independent replication attempt is positioned to make.

\section{Extended Analyses (Post-Review)}

The analyses in this section were added after an internal review round that identified two gaps: the primary result was reported only at a single view budget (N=50), and it did not rule out the possibility that the fixed mixed policy simply removes the labeled diagnostic content from a substantial fraction of augmented views. Except for the view-count scaling curve --- which was preregistered in the protocol as a secondary/descriptive analysis --- these analyses were specified and run after the final-test unsealing and are therefore reported as secondary / exploratory. Every number below traces to a machine- readable summary artifact (\texttt{artifacts/\allowbreak secondary\_\allowbreak analysis\_\allowbreak expansion/\allowbreak }, \texttt{artifacts/\allowbreak label\_\allowbreak preservation\_\allowbreak audit/\allowbreak }, \texttt{artifacts/\allowbreak component\_\allowbreak ablation/\allowbreak }) produced by a read-only script, and is checked by the same manuscript verification tool used for the Results section.

\begin{figure}[t]
\centering
\includegraphics[width=0.72\textwidth]{fig6_scaling_curve}
\caption{\textbf{View-count scaling curve (preregistered secondary).} Mean $\Delta$ accuracy vs.\ TTA view count $N$; the shaded band is the min--max over the 30 unmatched cells. Harm is largest at $N{=}1$, flattens by $N\approx25$, and never crosses zero.}
\label{fig:scaling}
\end{figure}

\textbf{View-count scaling curve.} Re-aggregating the already-computed 100-view sequence at N in {1, 2, 5, 10, 25, 50, 100} for the 30 unmatched-policy cells, naive mixed-policy mean-probability TTA reduced accuracy at every view count, in every cell, with each per-cell 95\% confidence interval excluding zero on the negative side at every N. The harm was largest at N=1 (mean -49.41 percentage points), decreased monotonically as views were added, and had flattened by roughly N=25; the preregistered primary condition N=50 (mean -42.56 percentage points) sits on this asymptote, and N=100 (mean -42.42 percentage points) is essentially identical. The curve never crosses zero and never reverses direction, so N=50 is close to the least-harmful operating point across the registered range rather than an unrepresentative choice (Figure 6).

\begin{figure}[t]
\centering
\includegraphics[width=0.72\textwidth]{fig8_label_preservation}
\caption{\textbf{Label-preservation audit (secondary / exploratory).} Human content-presence scores for 50 augmented views per dataset; the marker is the automated not-preserved rate, reported as a conservative upper bound the human check does not corroborate.}
\label{fig:labelpres}
\end{figure}

\textbf{Label-preservation audit.} For the frozen mixed policy at 28px we sampled 200 validation images per dataset, generated 16 augmented views each, recovered the actual sampled transform parameters, and applied a preregistered per-dataset "content-preserved" rule; a single annotator then scored a blinded, proxy-stratified subset of 50 augmented views per dataset on a three-point content-presence scale (not a diagnostic judgement). The labeled structure was judged clearly present in the large majority of views and judged lost ("gone") in 0.0\% of BloodMNIST views, 2.0\% of DermaMNIST views, and 4.0\% of PathMNIST views. The geometric transforms almost never displaced the labeled content: the random resized crop retained about 78\% of image area on average and close to 100\% of the image centre. Under the audit's preregistered interpretation rule --- "confound is material" if the human lost-content rate reaches 15\% or the automated not-preserved rate reaches 25\% --- all three datasets fell in the "unlikely to dominate" branch. The automated proxy flagged about 12.2\%, 12.5\%, and 13.4\% of views (BloodMNIST, PathMNIST, DermaMNIST) as not preserved, driven almost entirely by strong brightness/contrast jitter; the human check did not corroborate this as content loss, so the automated rate is reported only as a conservative upper bound (Figure 8).

\begin{figure}[t]
\centering
\includegraphics[width=0.72\textwidth]{fig7_component_decomposition}
\caption{\textbf{Per-augmentation-component decomposition (secondary / exploratory).} Geometric-only and intensity-only $\Delta$ accuracy at $N{=}50$ on the 12 Block A 28px cells, with the mixed-policy value marked. Both families harm accuracy on their own in all 12 cells.}
\label{fig:component}
\end{figure}

\textbf{Per-augmentation-component decomposition.} Splitting the mixed policy into its geometric family (flips, rotation, resized crop) and its intensity family (brightness/contrast jitter, Gaussian blur) and running each alone on the 12 Block A 28px cells (validation split, reusing the trained checkpoints), geometric-only TTA reduced accuracy by a mean of 26.20 percentage points and intensity-only TTA by a mean of 23.34 percentage points at N=50 (26.30 and 25.10 percentage points at the preregistered N=25); every one of the 12 cells showed a negative delta with a 95\% confidence interval excluding zero for both families at both view counts. Neither family dominated (their mean magnitudes differ by under 15\%), and the mixed-policy harm (mean -43.48 percentage points at N=50) was smaller than the sum of the two single-family harms, i.e. sub-additive (additivity residual +6.06 percentage points). Because the geometric-only family is content-preserving (per the audit above) yet still collapses accuracy by roughly a quarter, the harm is not attributable to the augmentation destroying labeled content; it is an input-distribution-shift effect --- these models, trained without augmentation, have little robustness to augmented inputs of either family, and naive averaging over such views amplifies rather than cancels that fragility (Figure 7).

\textbf{Other secondary conditions.} From the same already-computed final-test predictions we also extracted three source-study Appendix-B baselines and one robustness check, all preregistered as secondary and none a contribution of this work. BatchNorm-statistics adaptation reduced the harm relative to naive TTA (by roughly 11 to 27 percentage points on the two SmallCNN datasets) but did not restore accuracy (still roughly 10 to 16 percentage points below clean), so BatchNorm's running statistics are part but not all of the mechanism. Clean-image anchoring reduced the mean harm only modestly (to roughly -31 percentage points at N=50) and did not reproduce the large rescue the source study reported for that condition. Replacing mean-probability aggregation with majority vote or confidence-weighted averaging changed the mean harm by about one percentage point, so the effect is not an artifact of the aggregation rule. Naive TTA also roughly doubled expected calibration error and increased negative log-likelihood and multiclass Brier score by a factor of about two and a half; BatchNorm adaptation recovered much of the likelihood/Brier inflation but not the calibration-error inflation.

\textbf{What this resolves.} Together the label-preservation audit and the component decomposition answer the augmentation-severity question: the observed harm is not explained by the mixed policy removing the labeled class from augmented views. It is better described as these models having essentially no robustness to augmentation-induced distribution shift --- geometric or photometric --- which naive TTA then amplifies. This reframes the paper's account of the effect without altering the primary within-cell result, which is unchanged.

\section{Limitations}

\textbf{Design and coverage limitations.} Only three training seeds per cell; sample-level paired tests within a cell do not substitute for a model-seed population replication study, and our findings should not be read as bounding the variance across a hypothetical larger seed population. Coverage is limited to the specific MedMNIST subsets, architectures, and normalization/resolution combinations actually tested; a fixed augmentation policy (mixed: geometric + intensity) and a fixed TTA view budget (N=50) were used throughout, and we make no claim about other policies or view counts. The three secondary cross-condition comparisons (normalization, resolution, policy-matching) were specified as a fixed statistical procedure after validation-stage results were already observed, but before the official test split was opened; we disclose this ordering explicitly rather than presenting these comparisons with confirmatory status.

\textbf{Limitations of the post-review extension.} The per-augmentation- component decomposition was run on the validation split, not the test split, and covers only the 12 Block A 28px cells (SmallCNN, PathMNIST/BloodMNIST, three seeds) --- not ResNet-18, DermaMNIST, 64/128px, or the matched-policy arm --- so its family-level attribution should not be read as generalizing across architectures. The label-preservation audit used a single annotator, is limited to 28px, and its content- presence task is deliberately not a clinical judgement; its automated structural proxy over-flags on intensity jitter and is reported only as a conservative upper bound. Apart from the preregistered view-count scaling curve, every analysis in the "Extended Analyses" section was specified after the final-test unsealing and is exploratory, not confirmatory.

\textbf{Process incidents disclosed for completeness.} During execution, one final-test cell experienced an accidental final-test access incident on its first attempt (aborted, no scientific value persisted) and two further failed engineering attempts before a successful, audited completion; a shared-aggregation-contract correction was made mid-project after a defect was found, requiring a validation-metric-reconciliation step; all 39 canonical final-test results reported in this manuscript were produced under the final, corrected evaluator/aggregation pipeline, with the one cell computed under an earlier pipeline generation independently re-verified compatible via full recomputation rather than assumed compatible. No final-test scientific result was inspected by a human before the controlled, authorization-gated unsealing step that produced the canonical evidence package this manuscript reports from. A pre-existing test-suite gap (31 tests whose guard fixture incorrectly assumed the canonical outputs did not yet exist) was identified and fixed in the same project phase that produced this manuscript, as a test-only change with no effect on any scientific result. These incidents are recorded here, and in full detail in the project's audit trail, precisely so that engineering history does not have to be inferred or taken on faith; none of them altered which test-split data contributed to any reported number, and none required re-deriving any statistic.

\textbf{Scope of the novelty claim.} Our literature review, detailed in \texttt{docs/\allowbreak literature\_\allowbreak review.md} and re-verified for this manuscript in \texttt{paper/\allowbreak citation\_\allowbreak audit.md}, was not an exhaustive systematic search (e.g. no citation-graph traversal was performed); the novelty gap we identify (sec. "Related Work") is stated with that caveat and should not be read as a confirmed first.

\section{Reproducibility and Audit Trail}

This study's confirmatory statistical analysis, cross-condition secondary analysis, and scientific-result unsealing were each executed by a sealed, deterministic pipeline whose source code, environment, and inputs are bound together by cryptographic content fingerprints, verified unchanged before and after every execution step. The statistical analysis, cross-condition addendum, and unsealing steps each required an explicit, human-reviewed authorization artifact, verified before any result was computed or read; no scientific result in this manuscript was read from a partially-authorized or unverified pipeline state.

All numbers in the Results section above trace mechanically to the canonical scientific summary (\texttt{artifacts/\allowbreak final\_\allowbreak test\_\allowbreak scientific\_\allowbreak summary.json}) and the committed paper-evidence figures and tables (\texttt{artifacts/\allowbreak paper\_\allowbreak evidence/\allowbreak }), both of which are content-hash-bound in a machine-readable manifest (\texttt{artifacts/\allowbreak paper\_\allowbreak evidence/\allowbreak paper\_\allowbreak evidence\_\allowbreak manifest.json}). A read-only verification script (\texttt{paper/\allowbreak verify\_\allowbreak manuscript\_\allowbreak claims.py}, described in \texttt{paper/\allowbreak README.md}) cross-checks every numeric claim in this manuscript against that evidence package and rejects known-forbidden phrasings (e.g. "54 distinct cells," any secondary-significance language, any H4 claim); this manuscript was not considered complete until that script passed. The numeric claims in the "Extended Analyses" section are checked by the same script against their own read-only summary artifacts (\texttt{artifacts/\allowbreak secondary\_\allowbreak analysis\_\allowbreak expansion/\allowbreak summary.json}, \texttt{artifacts/\allowbreak label\_\allowbreak preservation\_\allowbreak audit/\allowbreak summary.json}, \texttt{artifacts/\allowbreak component\_\allowbreak ablation/\allowbreak summary.json}), each produced by a committed script that reads only already-sealed predictions, the validation split, or the trained checkpoints --- never the test-split model outputs beyond the already-unsealed canonical summary. No script or test described here accesses raw predictions, the underlying image datasets, or model checkpoints for the primary result --- only the already-sealed, already-verified summary and evidence-package artifacts.

\section{Conclusion}

Within a preregistered, statistically audited matrix of 30 distinct trained-model cells across three MedMNIST datasets, two normalization schemes, and two resolutions, naive mixed-policy test-time augmentation reduced test accuracy in every cell evaluated. A secondary, explicitly non-preregistered analysis suggests that matching the training-time augmentation policy to the test-time TTA policy mitigates this harm, while normalization and resolution show heterogeneous, dataset-dependent secondary patterns rather than a single consistent mitigation, and an internal positive-control replication did not reproduce the one previously reported case of TTA helping in this setting. A post-review extension establishes that the harm is stable across view counts from 1 to 100, that augmented views overwhelmingly retain the labeled structure, and that a content-preserving geometric-only augmentation family causes harm on its own --- so the effect is best understood as a lack of augmentation robustness in these models rather than as the augmentation destroying labeled content. We report all of this using an explicit evidentiary hierarchy that separates confirmatory within-cell results from secondary cross-condition comparisons, and we release the complete, hash-verified evidence and audit trail alongside this manuscript so that every reported number can be independently mechanically re-checked.

\section*{Supplementary Material Outline}

\emph{(Outline only in this phase, per \texttt{docs/\allowbreak phase2b\_\allowbreak manuscript\_\allowbreak claims\_\allowbreak and\_\allowbreak structure\_\allowbreak freeze.md} sec.6; full supplementary content is future work.)}

\begin{itemize}
  \item \textbf{Supplementary Table 1} (main-text Table 1 in full): experimental- design and evidence-classification table.
  \item \textbf{Supplementary Table 2}: complete 30-cell unmatched-policy table (dataset, resolution, normalization, seed, delta accuracy, 95\% CI, raw McNemar p, BH-adjusted p per member hypothesis family) --- \texttt{artifacts/\allowbreak paper\_\allowbreak evidence/\allowbreak tables/\allowbreak table\_\allowbreak 2\_\allowbreak unmatched\_\allowbreak policy.md}.
  \item \textbf{Supplementary Table 3}: complete matched-policy table (6 within-cell rows + 6 secondary DiD pairs) --- \texttt{artifacts/\allowbreak paper\_\allowbreak evidence/\allowbreak tables/\allowbreak table\_\allowbreak 3\_\allowbreak matched\_\allowbreak policy.md}.
  \item \textbf{Supplementary Table 4}: complete 12-pair normalization DiD table --- \texttt{artifacts/\allowbreak paper\_\allowbreak evidence/\allowbreak tables/\allowbreak table\_\allowbreak 4\_\allowbreak normalization.md}.
  \item \textbf{Supplementary Table 5}: complete 12-pair resolution DiD table --- \texttt{artifacts/\allowbreak paper\_\allowbreak evidence/\allowbreak tables/\allowbreak table\_\allowbreak 5\_\allowbreak resolution.md}.
  \item \textbf{Supplementary Table 6}: complete 3-seed BLOCK\_C table with the external +1.6pp descriptive reference footnote --- \texttt{artifacts/\allowbreak paper\_\allowbreak evidence/\allowbreak tables/\allowbreak table\_\allowbreak 6\_\allowbreak block\_\allowbreak c.md}.
  \item \textbf{Supplementary Table 7} (main-text Table 7 in full): claim adjudication table --- \texttt{artifacts/\allowbreak paper\_\allowbreak evidence/\allowbreak tables/\allowbreak table\_\allowbreak 7\_\allowbreak claim\_\allowbreak adjudication.md}.
  \item \textbf{Supplementary Note A}: full frozen training/evaluation hyperparameter table (\texttt{docs/\allowbreak phase2b\_\allowbreak protocol.md}).
  \item \textbf{Supplementary Note B}: full statistical analysis plan (\texttt{docs/\allowbreak statistical\_\allowbreak analysis\_\allowbreak plan.md}).
  \item \textbf{Supplementary Note C}: complete engineering and audit trail (fingerprint values, authorization chain, generation history) --- \texttt{docs/\allowbreak phase2b\_\allowbreak paper\_\allowbreak evidence\_\allowbreak package.md} and the referenced Phase 2B.6-2B.9 documents.
  \item \textbf{Supplementary Table 8}: view-count scaling curve, per-cell delta accuracy and 95\% CI at N in {1,2,5,10,25,50,100} --- \texttt{artifacts/\allowbreak secondary\_\allowbreak analysis\_\allowbreak expansion/\allowbreak tables/\allowbreak scaling\_\allowbreak curve.csv}.
  \item \textbf{Supplementary Table 9}: BatchNorm-adaptation, clean-anchoring, and aggregation-rule secondary conditions per cell --- \texttt{artifacts/\allowbreak secondary\_\allowbreak analysis\_\allowbreak expansion/\allowbreak tables/\allowbreak \{bn\_\allowbreak adaptation,anchored,aggregation\_\allowbreak ablation\}.csv}.
  \item \textbf{Supplementary Table 10}: label-preservation audit --- per-view proxy quantities (9,600 rows) and per-dataset human/automated rates --- \texttt{artifacts/\allowbreak label\_\allowbreak preservation\_\allowbreak audit/\allowbreak \{per\_\allowbreak view.csv,summary.json\}}.
  \item \textbf{Supplementary Table 11}: per-augmentation-component decomposition, 12 cells x {geometric, intensity} x {N=25, N=50} --- \texttt{artifacts/\allowbreak component\_\allowbreak ablation/\allowbreak per\_\allowbreak cell.csv}.
  \item \textbf{Supplementary Note D}: post-review extension protocols and findings --- \texttt{docs/\allowbreak phase2c\_\allowbreak secondary\_\allowbreak analysis\_\allowbreak expansion\_\allowbreak plan.md}, \texttt{docs/\allowbreak phase2c2\_\allowbreak label\_\allowbreak preservation\_\allowbreak audit\_\allowbreak protocol.md}, \texttt{docs/\allowbreak phase2c2\_\allowbreak component\_\allowbreak ablation\_\allowbreak addendum.md}, and the three corresponding \texttt{*\_\allowbreak findings.md}.
\end{itemize}
